\documentclass[11pt]{article}
\usepackage[margin=1.1in]{geometry}
\usepackage{amsmath,amssymb}
\usepackage{booktabs}
\usepackage[hidelinks]{hyperref}
\usepackage{parskip}

\newcommand{\code}[1]{\texttt{#1}}
\newcommand{\gf}{\mathrm{GF}(2)}

\title{\textbf{Deferred-Mask Nonlinear Gadgetization:\\ What We Tried and Why
It Failed}}
\author{local\_mixing project}
\date{July 21, 2026 \\ \small branch \code{ssg-gen-mix-clean} --- post-mortem
of the ``split-RG3'' attempt. Full design: \code{docs/NONLINEAR\_SHARE\_ENCODING.*}}

\begin{document}
\maketitle

\begin{abstract}
\noindent We tried to make the two-share ($\times2$) gadget hide the
degree-1 computational-progress diagonal by encoding each logical value with
a \emph{deferred nonlinear mask} --- the RG3 randomizer split in time. The
construction is correct and demonstrably hides values \emph{between} their
uses, but on the $n{=}16$ degree-1 reconstruction grid the progress diagonal
\textbf{survives fully} (rho $=1.00$, ridge depth unchanged), and it is
independent of both masking coverage and stacking depth. The cause is the
\emph{peek}: a masked value must be momentarily reconstructed at each gate
that uses it as a control, and the reconstruction diagonal is precisely the
locus of those uses. The lesson is that one cannot hide a value at the
instant it is consumed if it is consumed by reconstructing it; removing the
diagonal requires computing on the masked operand \emph{without}
reconstructing it.
\end{abstract}

\section{The problem it was meant to solve}

The gadget carries every logical value $v$ as an XOR pair $v = c_0 \oplus
c_1$. That makes $v$ a \textbf{degree-1 function of the physical wires at
every snapshot}, no matter how nonlinear the transitions between snapshots
are. The measured consequence is the computational-progress \emph{diagonal}
in the \code{hmap\_affine} reconstruction maps --- a low-error valley whose
position advances with the original circuit's prefix --- and it survives
millions of \code{fmix} mixing moves with alignment intact (rho $= 1.00$).
Nonlinear re-randomizers (RGs) are provably orthogonal to this: they make
\emph{crossing} the body expensive for a low-degree adversary, but reading a
\emph{snapshot} crosses nothing. The leak is in the \textbf{encoding}, not
the transitions, so the fix has to change the encoding.

\section{The idea: defer the RG3 randomizer in time}

The nonlinear RG3 already XORs a degree-2 randomizer $r_1 \vee \bar r_2$ into
\emph{both} carriers of a value --- a refresh that leaves the value unchanged
(the two writes cancel in the XOR), which is exactly why RG3 re-randomizes
the shares but cannot mask the value.

The move (``RG4'', or ``split-RG3''): emit those two writes \textbf{separated
in time}. Inject the term into \emph{one} carrier now; compensate on a
carrier later. Between injection and compensation the value reads as $v
\oplus (r_1 \vee \bar r_2)$ --- \textbf{degree 2} in the wires, with random
support. If at every snapshot most values carry such a pending term, the
whole state is degree-2 rather than degree-1, and the affine adversary that
draws the diagonal should lose it.

Three refinements made the move workable (details in the design note):

\begin{itemize}
\item \textbf{Value-sourced.} The randomizer reads two \emph{logical source
values} $v_{r1} \vee \lnot v_{r2}$, rebuilt from their current carriers at
each emit. Because RG1/RG2/RG3 preserve logical values, the mask term is
invariant under the entire RG stream; only a computation gate (CG) that
\emph{recomputes} a source disturbs it, and that schedule is known in
advance.
\item \textbf{Lookahead.} Sources are chosen to be the values whose next
recomputation is farthest out (finished outputs are ideal), so a mask
survives from injection to a distant compensation instead of collapsing back
into a plain RG3.
\item \textbf{Peek on reads.} When the virtual circuit \emph{uses} $v$ as a
control, the gate needs the true $v$: the mask is momentarily removed
(un-mask), the ordinary gate runs on the clean value, and the same term is
re-applied (re-mask). The mask keeps its identity and lifetime.
\end{itemize}

The knob $k$ stacks $k$ independent degree-2 terms on a value, intended to
drive the best affine reconstruction error toward $\tfrac12$ by piling-up
($k{=}1 \to 0.25$, $k{=}3 \to 0.44$).

\section{What actually worked}

The machinery is correct and does part of what it should:

\begin{itemize}
\item \textbf{Function-preserving.} Every injection is compensated before
decode, so the gadget computes the identical function; the pinned zero-slice
contract check passes on masked builds and the test suite is green
($161/161$); vocabulary is clean (no bare X gates).
\item \textbf{It hides values \emph{between} their uses.} On $n{=}16$ the
masking raises the overall reconstruction error toward hidden --- mean map
error $H$ rises from $0.432$ (off) to ${\sim}0.449$, and the diagonal's
contrast against the background drops from $2.38\sigma$ to ${\sim}1.7\sigma$.
The \textbf{off-diagonal background gets more uniformly hidden.}
\end{itemize}

So the idea is not incoherent and it is not a bug: it demonstrably masks. The
question was whether it masks the thing that matters --- the diagonal.

\section{The experiment}

A fixed original circuit $C$ (160 random g57 gates, $n{=}16$) was gadgetized
with masking \textbf{off} and \textbf{on} across a grid of coverage and
stacking, then each gadget $G$ was compared to $C$ with \code{hmap\_affine}
at degree 1 and read with the ridge measure (\code{depth} $=$ per-row
prominence of the diagonal, \code{rho} $=$ its monotonic alignment,
\code{perm z} $=$ significance vs.\ a column-shuffle null, \code{contrast} $=$
diagonal-vs-background separation). The ridge measure --- not the mean --- is
the agreed way to read these maps.

\begin{center}
\begin{tabular}{lccccc}
\toprule
config & mean $H$ & depth & rho & perm $z$ & contrast \\
\midrule
\textbf{off}       & 0.432 & 0.348 & 1.00 & 4.50 & 2.38 \\
cov 0.5,\ \ $k{=}1$ & 0.443 & 0.316 & 1.00 & 4.54 & 1.89 \\
cov 0.75, $k{=}1$  & 0.449 & 0.316 & 1.00 & 4.49 & 1.64 \\
cov 1.0,\ \ $k{=}1$ & 0.445 & 0.357 & 1.00 & 4.53 & 1.84 \\
cov 0.75, $k{=}2$  & 0.441 & 0.344 & 1.00 & 4.53 & 2.05 \\
cov 0.75, $k{=}3$  & 0.449 & 0.347 & 1.00 & 4.48 & 1.92 \\
\bottomrule
\end{tabular}
\end{center}

\section{The result: the diagonal survives}

\textbf{The masking does not break the degree-1 diagonal.} \code{rho} stays
$1.00$ and \code{depth} stays ${\sim}0.32$--$0.36$ in every configuration,
statistically indistinguishable from the unmasked gadget. The diagonal
lightens slightly (deep-red ${\sim}0.15$ toward orange ${\sim}0.25$) but
never dissolves.

Two facts make this a trustworthy negative rather than a tuning artifact:

\begin{itemize}
\item \textbf{Coverage-independent.} Raising the masked fraction from $0.5$ to
$1.0$ does not reduce the depth. (Aside: \code{--mask-cov 1.0} only reaches
${\sim}0.75$ \emph{actual} coverage --- value-sourcing structurally needs an
unmasked pool to draw sources from --- but since the diagonal does not track
coverage at all, that limitation is not what carries it.)
\item \textbf{Stacking-independent.} Going from $k{=}1$ to $k{=}3$ does not
reduce the depth, even though piling-up predicts $k{=}3$ should push masked
values to ${\sim}0.44$ error and flatten the diagonal.
\end{itemize}

The verdict is a clean negative: the deferred value-sourced mask hides the
background but leaves the progress diagonal fully legible, and no amount of
\emph{more masking} changes that.

\section{Why it failed: the peek carries the diagonal}

Put the three observations together --- the background is hidden, the diagonal
survives, and it is immune to both coverage and stacking. That combination
points at one mechanism.

The reconstruction diagonal \textbf{is the locus of operand-uses}: the point
in $G$ that aligns with step $i$ of $C$ is the gate that computes step $i$,
and that gate \textbf{peeks its operands} --- un-masks them so the ordinary
gate can read their true values. Those exposed operand bits make $C_i$
affinely recoverable at exactly the aligned $G_j$, regardless of how many
\emph{other} values are masked (coverage) or how many terms sit on the target
value (stacking). The diagonal rides on the peeked operands, not on the
masked targets, so the two knobs that deepen target-masking cannot touch it.

Stated plainly: \textbf{one cannot hide a value at the instant it is consumed
if it is consumed by reconstructing it.} The peek was the pragmatic shortcut
that made the CG variant-agnostic and simple; the experiment shows it is
precisely the wrong trade at the one place --- the use-point --- that the
diagonal measures.

This closes the loop with an earlier finding: the diagonal is a positional /
linear-sharing leak, and the \emph{legacy} g57 gadget hides it at degree 2
not by luck but because it computes on the \textbf{shares} --- every gate
reads at most one carrier of each value and never reconstructs an operand
(gate-local non-completeness). That is the whole reason sharing-based gadgets
compute on shares in the first place. The deferred mask tried to sidestep
share-native computation with a momentary reconstruction, and the peek is
where the sidestep breaks.

\section{The lesson and what it points to}

The deferred-mask encoding is sound engineering and it does mask between
uses; it is not the mechanism that removes the diagonal. Removing the
diagonal requires computing on a masked operand \textbf{without reconstructing
it}:

\begin{itemize}
\item a \textbf{mask-aware CG} that folds the mask correction ($\mu \cdot C$
for a masked operand) into the gate's algebra, keeping the operand masked
throughout; or
\item a \textbf{share-native CG} in the style of the legacy g57 gadget, whose
gate-local non-completeness is exactly what buys degree-2 hiding.
\end{itemize}

Either is a larger change than the peek. The value of this attempt is that it
turned an intuition (``the peek exposure is a minor localized cost'') into a
measured, falsifiable conclusion: the peek is the wall, and the next
prototype should keep the operands masked through the gate and re-run this
same grid --- the diagonal depth should drop if and only if that is done.

\end{document}
